\chapter{Backup, logs and troubleshooting}\label{ch:backup}

\emph{You want to be able to rebuild Agern after a disaster, and to find out quickly what is wrong
when something does not work.}

\section{Back up}\label{sec:backup}

Back up two things, and keep them apart (plus the project files, below, which go with the database).

\begin{enumerate}
  \item \textbf{The database}, for example every night (as a cron job or systemd timer):
\begin{lstlisting}
sudo -u postgres pg_dump -Fc factory > /var/backups/agern-db-$(date +%F).dump
\end{lstlisting}
  \item \textbf{The configuration} \file{/etc/agern/agern.env}, above all \envvar{FACTORY_MASTER_KEY}.
    Every stored Claude credential, GitHub token and two-factor secret in the database is encrypted
    with it. A database backup without the key is useless, and the key together with the database
    gives access to everybody's credentials: keep the key in a password manager or secrets store, not
    next to the dumps.
\end{enumerate}

One part of the data directory must be backed up as well: \file{/var/lib/agern/project-files}, the
files uploaded on the projects' \ui{Files} pages (section~\ref{sec:project-files}). They exist
nowhere else and no agent can recreate them, and their names and hashes are in the database while
their bytes are only in this directory, so copy it right after the database dump and restore the
two together (\cmd{rsync -aH} or \cmd{tar} keep its hard links). Plan disk space for them: by
default a project may store 5~GB and all projects together 20~GB, as set by
\envvar{FACTORY_PROJECT_FILES_MAX_FILE_MB}, \envvar{FACTORY_PROJECT_FILES_MAX_PROJECT_MB} and
\envvar{FACTORY_PROJECT_FILES_MAX_TOTAL_MB}.

The rest of the data directory \file{/var/lib/agern} holds checkouts and caches that agents
recreate. Include it only if you want interrupted chats to resume after a restore. Snapshots of the whole virtual
machine are a good addition to, not a replacement for, the database dump.

\section{Restore on a new machine}\label{sec:restore}

\begin{enumerate}
  \item Install PostgreSQL and the Agern package as in chapter~\ref{ch:install}, but do not start Agern.
  \item Put your saved \file{/etc/agern/agern.env} back, owned by \texttt{root:agern} with mode 0640:
\begin{lstlisting}
sudo install -o root -g agern -m 0640 agern.env /etc/agern/agern.env
\end{lstlisting}
  \item Create the role (with the password from that file) and an empty database, restore the dump
    and start:
\begin{lstlisting}
sudo agern setup-db
sudo -u postgres pg_restore --no-owner --role=factory -d factory agern-db-2026-11-24.dump
sudo agern start
\end{lstlisting}
  \item Install the agents' tools again (section~\ref{sec:agent-tools}) and run \cmd{sudo agern doctor}.
\end{enumerate}

\section{Logs and the doctor}\label{sec:logs}

\begin{enumerate}
  \item Agern logs to the system journal. Show the latest lines, or follow the log live:
\begin{lstlisting}
agern logs -n 200
agern logs -f
\end{lstlisting}
    (\cmd{journalctl -u agern} shows the same.) Automatic upgrades log to
    \cmd{journalctl -u agern-upgrade}, sandboxes to \cmd{journalctl -u agern-sandbox}.
  \item When anything is odd, ask the doctor first. It checks the configuration, the database login,
    Java, \texttt{claude}, \texttt{git} and \texttt{gh}, the headless browser, the data directory,
    the \texttt{agern} user's privileges, sandboxes and whether the server answers:
\begin{lstlisting}
sudo agern doctor
\end{lstlisting}
  \item \cmd{agern status} says whether the service runs and whether its work is in progress.
\end{enumerate}

\section{Troubleshooting}\label{sec:troubleshooting}

The messages below are exactly what Agern shows.

\begingroup
\setlength{\parskip}{0pt}
\small
\renewcommand{\arraystretch}{1.3}
\begin{longtable}{@{}>{\raggedright\arraybackslash}p{0.33\linewidth}>{\raggedright\arraybackslash}p{0.25\linewidth}>{\raggedright\arraybackslash}p{0.36\linewidth}@{}}
  \toprule
  What you see & Cause & Fix\\
  \midrule
  \endhead
  \msg{The Agern server is not answering.} (web UI, with \ui{Try again}) & The server is stopped,
  starting, or crashed. & \cmd{agern status}; \cmd{sudo agern start}; read \cmd{agern logs}.\\
  \msg{Origin not allowed} when signing in or saving anything & The address in the browser differs
  from \envvar{FACTORY_PUBLIC_URL}. & Open exactly that address, or correct the variable and
  \cmd{sudo agern restart} (section~\ref{sec:public-url}).\\
  \texttt{agern start} fails; the log says \msg{Agern: configuration error: Environment variable
  FACTORY\_GITHUB\_ORG is required} & A required setting is missing or invalid; Agern exits instead
  of restarting in a loop. & Fix the named variable in \file{/etc/agern/agern.env}; run
  \cmd{sudo agern doctor}.\\
  \msg{Wrong username or password} although the password is right, and the doctor warns
  \msg{FACTORY\_AUTH has 'os' but agern is not in group 'shadow'; PAM cannot check other users'
  passwords (see the install docs)} & PAM cannot read \file{/etc/shadow} for \texttt{agern}. &
  \cmd{sudo usermod -aG shadow agern} and \cmd{sudo agern restart}, or use LDAP
  (section~\ref{sec:first-admin}).\\
  \msg{Too many wrong codes. Two-factor sign-in is locked for 15 minutes.} & Five wrong two-factor
  codes in a row. & Wait 15 minutes; check the phone's clock; use a recovery code; or ask an admin
  for \ui{Reset 2FA} (section~\ref{sec:users}).\\
  \msg{GitHub rejected the token (401)} in Settings & The token is wrong, expired or revoked. &
  Create a new token (section~\ref{sec:github-token}) and paste it.\\
  \ui{No repositories visible} on the projects page & Your token cannot read any repository of the
  organization (fine-grained token without access to it, or the organization restricts tokens). &
  Give the token access to the organization's repositories; ask an organization owner to approve it.\\
  \msg{Register a GitHub token in Settings first} & You have no GitHub token registered. & Add it
  under \ui{Settings} $\rightarrow$ \ui{GitHub}.\\
  \msg{Register a Claude credential in Settings first} when starting a run or release, or
  \ui{Chat is disabled} & You have no Claude credential registered. & Add it under \ui{Settings}
  $\rightarrow$ \ui{Claude} (section~\ref{sec:claude-credential}).\\
  \msg{You already have 8 agents running} & You reached \envvar{FACTORY_MAX_AGENTS_PER_USER}. &
  Close chats you no longer need, wait for runs to finish, or raise the limit.\\
  The doctor says \msg{claude not found on \ldots} & Claude Code is not installed for the
  \texttt{agern} user. & Run the command the doctor prints (section~\ref{sec:agent-tools}).\\
  The doctor warns \msg{FACTORY\_PUBLIC\_URL is http://localhost:8080 --- set the URL people
  type in their browser, or every browser request from the LAN is rejected} & The public URL was
  never set. & Set \envvar{FACTORY_PUBLIC_URL} (section~\ref{sec:minimal-config}).\\
  \msg{agern: error: not restarting now, the server is busy: \ldots{} Try again later, or use
  -{}-force (interrupted runs fail)} & \cmd{agern upgrade} or \cmd{agern sandbox disable} while
  work is in progress. & Wait until runs finish, or add \cmd{--force} (chapter~\ref{ch:upgrade}).\\
  A run page says its sandbox definition is not approved & \file{.agern/sandbox.yml} is new or
  changed. & An admin approves it on the project's \ui{Sandbox} page (section~\ref{sec:sandbox-approve}).\\
  \bottomrule
\end{longtable}
\endgroup

When none of this helps,
\repodoc{INSTALL-LINUX.md}{\texttt{INSTALL-LINUX.md}} has more detail on every part of the
installation, and \repodoc{ARCHITECTURE.md}{\texttt{ARCHITECTURE.md}} on how Agern works.

\whatyouhave
Your data is backed up and you know how to restore it, where the logs are, and how to fix the
common problems. You have read the whole path: Agern is installed, configured and in use. The
appendices are for looking things up: the glossary (appendix~\ref{app:glossary}) and the
configuration quick reference (appendix~\ref{app:config}).
